%==============================================================================
% Section 2: Lorentzian-native setup and topology conventions
%==============================================================================
\section{Lorentzian-native setup and topology conventions}
\label{sec:setup}

\subsection{Reduced Lorentzian frame}
\label{sec:setup:frame}

We use the orthonormal frame of Lorentzian signature $(-,+,+,+)$.  The
coframe is written as
\begin{equation}
  e^{0} = N(t)\,dt, \qquad e^{i} = q(t)\,\sigma^{i}_{\topo},
\end{equation}
where $N(t)$ is the lapse, $q(t)$ is the isotropic scale common to all
topologies $\topo$, and $\sigma^{i}_{\topo}$ are the left-invariant coframes
on $\topo$.  Torsion variables are taken to be the axial component $\eta(t)$
and the vector-trace component $V(t)$, giving the pure branches $\AX$ and
$\VT$ and the mixed branch $\MX$.

The normalisation of the spatial 3-curvature is fixed by
\begin{equation}
  {}^{(3)}R = \frac{6\chi}{q^{2}},
\end{equation}
and $\chi$ is used as the topology scalar throughout this paper.  The scalar
$\chi$ is derived from the $\EH$ (pure Einstein--Hilbert) sector only and
does not depend on the torsion variables.

\subsection{Topology structure constants and $\chi$ values}
\label{sec:setup:chi}

The structure constants of the four topologies are specialisations of the
five-parameter LZ-native isotropic family, following the class-A / class-B
classification of Bianchi-type homogeneous spaces~\cite{EllisMacCallum1969}
and the three-topology mode dictionary developed earlier in the DPPU
series~\cite{paper3ec}:
\begin{equation}
  C^{1}{}_{23}=\frac{a}{q},\quad
  C^{2}{}_{31}=\frac{b}{q},\quad
  C^{3}{}_{12}=\frac{c}{q},\quad
  C^{2}{}_{12}=\frac{u}{q},\quad
  C^{3}{}_{13}=\frac{v}{q},
\end{equation}
where the subscripts are antisymmetric.  Then
\begin{equation}
  {}^{(3)}R\,q^{2}
  = -\tfrac{1}{2}\!\left(a^{2}-2ab-2ac+b^{2}-2bc+c^{2}+4u^{2}+4uv+4v^{2}\right)
  =: -\tfrac{1}{2}\,\Disc,
\end{equation}
and comparison with the defining relation ${}^{(3)}R=6\chi/q^{2}$ gives
\begin{equation}
  6\chi = -\frac{\Disc}{2},
  \qquad
  \chi = -\frac{\Disc}{12}.
\end{equation}
The specialised values for each topology are listed in
Table~\ref{tab:topology_chi}.

\begin{table}[H]
\centering
\begingroup
\renewcommand{\arraystretch}{2.0}
\begin{tabularx}{\linewidth}{lcrYY}
\toprule
Topology & $(a,b,c,u,v)$ & $\chi$ & Reduction status & Caveat \\
\midrule
$\Sthree$   & $(4,4,4,0,0)$  & $4$    & closed isotropic reference & --- \\
$\Tthree$   & $(0,0,0,0,0)$  & $0$    & flat reference             & --- \\
$\Nilthree$ & $(0,0,-1,0,0)$ & $-1/12$ & isotropic-scale local branch & anisotropic EOM left to future work \\
$\Solthree$ & $(0,0,0,1,-1)$ & $-1/3$  & isotropic-scale local branch & compact quotient outside the scope of this paper \\
\bottomrule
\end{tabularx}
\endgroup
\caption{Topology conventions, $\chi$ values, and caveats.}
\label{tab:topology_chi}
\end{table}

We treat all four geometries in a common reduced-topology table.  For
$\Nilthree$ and $\Solthree$, the entries refer to the isotropic-scale
reductions used in this paper; global quotient and anisotropic extensions
are left to future work.

\subsection{EC+NY reduced action}
\label{sec:setup:action}

The Lorentzian EC+NY reduced action~\cite{Hehl1976,NiehYan1982} for each
topology $\topo$ and torsion mode is written in the form
\begin{equation}
  S_{\topo} = \int dt\, A_{\topo}\, q\!\left(N\,F_{\topo}(q,\eta,V;\chi)
  - \frac{\dot q^{2}}{N}\right).
\end{equation}
Here $A_{\topo}$ is a spatial normalisation factor independent of $\chi$.
The effective potential $F_{\topo}$ is determined by the topology and the
torsion branch, and its explicit form is used in
Section~\ref{sec:hamiltonian}.  The auxiliary variables $\eta(t)$ and $V(t)$
carry no velocity terms in the action and are treated algebraically as an
auxiliary sector.

The Nieh--Yan term $\thetaNY$~\cite{NiehYan1982} does not appear in the bulk
$F_{\topo}$ on $\AX$ and $\VT$, and appears on $\MX$ only as a mixing term
$\propto V\eta\,\kap^{2}\thetaNY$.  In this paper the bulk dynamical
analysis is taken as primary, and $\QNY$ is separated as an endpoint /
boundary channel (Section~\ref{sec:pchannel:QNY}).
